% ===== 原文 pp.32–34 =====
% =====================================================================
%  tr05a —— 原文 pp.32–34（§4.2 后半：4.2.2–4.2.3 起）
%  编号覆盖：式 (4.10)、式 (4.11)、定理 4.12、式 (4.13)
%  起点：p32 第 2 句 "Having mentioned the spectral data…"
%        （p32 首句 "…quotient spectral curves, as shown in [89,91]." 由 tr-04 补完，本块不重复）
%  终点：p34 末段；p35 由 tr05b.tex 接
% =====================================================================

% [原文 p.32]
前几章已经提及这些对象的谱数据，本节余下的篇幅将用于讨论下文所述的第三种可能性。
% [Tutor Note: p.32 第 3 行原文 "we shall dedicate the reminder of this section" 中 "reminder"
%  应为 "remainder"（本节余下的部分），属形近笔误；译文按正意（"本节余下的篇幅"）译出，正文不擅改。]

\begin{itemize}[leftmargin=1.4em,itemsep=1pt,topsep=1pt]
\item \textbf{与正则纤维的交为空。} 由实 Higgs 丛的结构（见 [89]），对于除 $U(p,p)$ 与 $\SU(p,p)$ 之外的实形式，非分裂（non-split）Higgs 丛的特征多项式将给出可约的谱曲线，因而其 $(\mathrm{B},\mathrm{A},\mathrm{A})$-膜完全落在 Hitchin 纤维化的判别式轨迹（discriminant locus）之上。可以证明，这些膜与 Hitchin 纤维化的非正则纤维以不同方式相交；下面我们考虑三种不同的情形：
  \begin{itemize}[label=--,leftmargin=1.8em,itemsep=1pt,topsep=1pt]
  \item \textbf{阿贝尔谱数据。} 这类数据的例子来自拟分裂但非分裂（quasi-split but not split）的实形式。对于这些群，[82] 表明：群 $\SU(p,p+1)$ 与 $\SO(p,p+2)$ 的 cameral 数据应当是阿贝尔的；而 $\SO(p,p+2)$-Higgs 丛的阿贝尔谱数据在 [11] 中有描述。

  \item \textbf{非阿贝尔谱数据。} 某些群的 Higgs 丛会携带自然的非阿贝尔谱数据，用以描述其 $(\mathrm{B},\mathrm{A},\mathrm{A})$-膜与 Hitchin 纤维化的非正则纤维之交。为研究这一情形，可以考察一般特征空间为 $r$ 维、且对某个多项式 $p(x)$ 满足 $\det(xI-\Phi) = p(x)^r$ 的 Higgs 丛 $(E,\Phi)$：由于 Hitchin 映射是满射，这类情形必定出现。当 $r=2$ 时，通过推广 [89] 中的方法，[64] 表明：若考虑 $\Sigma$ 上对应于和乐落在实李群 $G = Sp(2p,2p)$、$\Sl(p,\HH)$ 与 $\SO(2p,\HH)$（后两者又分别记作 $SU^*(2p)$ 与 $SO^*(2p)$）中的平坦联络的 Higgs 丛，则会出现一种“非阿贝尔化”（nonabelianization）程序，其谱数据由某条谱曲线上的秩 2 Higgs 丛的某个子空间给出。

  \item \textbf{阿贝尔与非阿贝尔谱数据。} 有趣的是，某些 $(\mathrm{B},\mathrm{A},\mathrm{A})$-膜与 Hitchin 纤维化的一般交需要同时用阿贝尔数据与非阿贝尔数据来描述。特别地，这正是 [11] 中所描述的 $\SO(p+q,p)$-Higgs 丛与 $\Sp(2p+2q,2p)$-Higgs 丛的情形。在此情形下，Higgs 场一般有一个 $2q$ 维的核，其特征多项式含有一个次数为 $4p$ 的多项式，而该多项式一般定义一条光滑曲线。通过考虑这条光滑曲线及其上的阿贝尔数据，再考虑一条带非阿贝尔数据的辅助谱覆盖，[11] 用“Cayley 与 Langlands 型对应”（Cayley and Langlands type correspondences）给出了该膜的几何描述，下文将对此加以说明。
  \end{itemize}
\end{itemize}

下面我们将给出一些完全落在 Hitchin 纤维化判别式轨迹之上的膜的例子，主要遵循 [64] 与 [11] 的工作。

% [原文 p.33]
\para{4.2.2.\ 非阿贝尔化} 按照定义 2.13，考虑以下情形：

\begin{itemize}[leftmargin=1.4em,itemsep=1pt,topsep=1pt]
\item $\Sl(n,\HH)$-Higgs 丛：由一个秩 $2n$ 的辛向量丛 $(V,\omega)$ 与一个 Higgs 场构成，后者对辛转置满足 $\Phi = \Phi^T$。用 $\omega$ 把 $V$ 与 $V^*$ 等同起来，这意味着 $\Phi = \omega^{-1}\varphi$，其中 $\varphi \in H^0(\Sigma, \Lambda^2V\otimes K)$。

\item $\SO(2p,\HH)$-Higgs 丛：其向量丛为 $V = W\oplus W^*$，其中 $W$ 是秩 $2p$ 的向量丛。内积由 $W$ 与 $W^*$ 之间的自然配对定义。Higgs 场具有形式 $\Phi(w,\xi) = (\beta(\xi),\gamma(w))$，其中 $\beta : W^*\to W\otimes K$ 与 $\gamma : W\to W^*\otimes K$ 均为斜对称的。

\item $\Sp(2p,2p)$-Higgs 丛：其向量丛为 $W_1\oplus W_2$，其中 $(W_1,\omega_1)$、$(W_2,\omega_2)$ 是辛秩 $2p$ 的向量丛；Higgs 场具有形式 $\Phi(v,w) = (\beta(w),\gamma(v))$，其中 $\beta : W_2\to W_1\otimes K$、$\gamma : W_1\to W_2\otimes K$，并且用辛转置有 $\beta = -\gamma^T$。
\end{itemize}

上述 Higgs 丛的谱数据在 [64] 中用秩 2 向量丛加以描述，并可概括如下：考虑亏格 $g \geqslant 2$ 的紧黎曼面 $\Sigma$ 上秩为 $2n$ 的 $\Sl(2n,\C)$-Higgs 丛 $(V,\Phi)$，其 Higgs 场的特征多项式为
% [Tutor Note: p.33 原文 "considering rank an SL(2n,C)-Higgs bundles" 中 "rank an" 为衍文/漏字，
%  据下文 "by letting n = 2p" 及 $\Sl(2n,\C)$-丛的秩即 $2n$ 判断，原意应为 "a rank 2n"；译文按正意写出。]
\begin{equation}
    \det(\Phi-\eta) = (\eta^n + a_2x^{n-2} + \cdots + a_n)^2. \tag{4.10}
\end{equation}
% [Tutor Note: p.33 式 (4.10) 的 $a_2x^{n-2}$（以及 p.34 式 (4.13) 与式 (4.13) 上方无编号公式中的
%  $a_1x^{2p-2}$）中出现的 $x$ 在全文从未定义；谱曲线的变量是 Tot$(K)$ 上的 重言截面 $\eta$
%  ——对照 p.21 的式 (3.8)，同一个特征多项式在那里写作
%  $\det(\eta-\Phi)=\eta^q(\eta^{2p}+a_2\eta^{2p-2}+\cdots+a_{2p})$，可见 $x$ 应为 $\eta$。
%  这里照原文排 $x$，请读者注意。]

应当指出，群 $\Sl(n,\HH)$、$\SO(2p,\HH)$ 与 $\Sp(2p,2p)$ 的 Higgs 丛都处在这一框架之内（取 $n = 2p$ 即可）。令 $a_i \in H^0(\Sigma, K^i)$、$\eta$ 为与之前相同的 重言截面，则上述多项式一般定义一条光滑谱曲线
% [Tutor Note: p.33 "lie within tis setting" 中 "tis" 为 "this" 之误（漏字），正文按正意译。]
\begin{equation}
    S = \{\sqrt{\det(\Phi-\eta)} = 0\}. \tag{4.11}
\end{equation}
% [Tutor Note: 式 (4.11) 的根号写法值得留意：由 (4.10)，$\det(\Phi-\eta)$ 是平方 $(\eta^n+\cdots+a_n)^2$，
%  故 $\sqrt{\det(\Phi-\eta)}=0$ 与 $\eta^n+a_2x^{n-2}+\cdots+a_n=0$ 定义同一条（既约）谱曲线；
%  原文保留根号记号，这里照排。]

\begin{thm}{4.12（[64]）} 考虑特征多项式如 (4.10) 所示、秩为 $2n$ 的 $\Sl(2n,\C)$-Higgs 丛 $(E,\Phi)$。给定由 (4.11) 定义的曲线 $S$ 上秩 2 的向量丛 $E$，则 $\eta : E \to E\otimes\pi^*K$ 的直像定义 $\Sigma$ 上的半稳定 $\Sl(n,\HH)$-Higgs 丛，当且仅当
\begin{enumerate}[label=(\roman*),leftmargin=2.2em]
\item $\Lambda^2E \cong \pi^*K^{n-1}$；
\item $E$ 是半稳定的。
\end{enumerate}
当 $n = 2p$ 时，(4.11) 中的曲线 $S$ 带有自然对合 $\sigma : \eta \mapsto -\eta$。此时，$\eta : E \to E\otimes\pi^*K$ 的直像定义 $\Sigma$ 上的半稳定 $\SO(2p,\HH)$-Higgs 丛，当且仅当 (i) 与 (ii) 成立，并且
\begin{enumerate}[label=(\roman*),leftmargin=2.2em,start=3]
\item $\sigma^*E \cong E$，其中在 $\Lambda^2E = \pi^*K^{4p-1}$ 上的诱导作用为平凡的。
\end{enumerate}
最后，$\eta : E \to E\otimes\pi^*K$ 的直像定义半稳定的 $\Sp(2p,2p)$-Higgs 丛，当且仅当 (i) 与 (ii) 成立，并且
\begin{enumerate}[label=(\roman*),leftmargin=2.2em,start=4]
\item $\sigma^*E \cong E$，其中在 $\Lambda^2E = \pi^*K^{2m-1}$ 上的诱导作用为 $-1$。
\end{enumerate}
\end{thm}
% [Tutor Note: p.33 定理 4.12 的 (iv) 中指数 $2m-1$ 里的 $m$ 在全文未出现；由 (iii) 的 $4p-1$
%  （$n=2p$ 时即 $2n-1$）推测此处应为 $2n-1$（或 $2p-1$，视 $\Sp(2p,2p)$ 丛的秩而定）。
%  照原文排 $m$，正确形式请对照 [64] 核定。]

有趣的是，正如 [89] 与 [64] 所指出的，在考虑 $\Sp(2p,2p)$-Higgs 丛时，$\sigma$ 在 $\Lambda^2E$ 上以 $-1$ 作用，于是在不动点处有互不相同的 $+1$ 与 $-1$ 特征空间。按照 [3]，这定义了
% [原文 p.34]
商曲线 $\pi' : \bar S := S/\sigma \to S$ 上的一个秩 2 丛，该商曲线在不动点处带有抛物结构：其旗由 $-1$ 特征空间给出，抛物权为 $1/2$；因此这与第 1 章所述的抛物 Higgs 丛密切相关。商曲线上 $E$ 的直像给出一个秩 2 向量丛，它可按 $\sigma$ 的诱导作用下的不变截面与反不变截面分解为线丛 $E^+ \oplus E^-$。秩 $4p$ 的向量丛可以通过商曲线还原：取 $\bar\pi(E^+) = W_1$ 与 $\bar\pi(E^+) = W_2$，其中 $\bar\pi : \bar S \to \Sigma$ 是秩 $2p$ 的投影。此外，由 Lefschetz 不动点公式，$W_1$、$W_2$ 的次数可以用 $E^\pm$ 的次数表示。
% [Tutor Note: p.34 第一行作 $\pi' : \bar S := S/\sigma \to S$，但商曲线应投影到 $\Sigma$；同段下文
%  又写 $\bar\pi : \bar S\to\Sigma$，故箭头后的 $S$ 疑为 $\Sigma$ 之误（漏印 Sigma）。照原文排。]
% [Tutor Note: p.34 "rank 4p vector bundle can be recovered" 缺主语冠词 "A"（应为 "A rank 4p…"）；
%  又 "taking $\bar\pi(E^+)=W_1$ and $\bar\pi(E^+)=W_2$" 中第二个 $E^+$ 与上句 $E^+\oplus E^-$ 的分解不符，
%  按 $W_1$、$W_2$ 分别来自不变与反不变部分判断应为 $\bar\pi(E^-)=W_2$。正文照原文排 $E^+$，在此更正。]

\para{4.2.3.\ Cayley 与 Langlands 型对应} 回忆一下，$\SO(p+q,p)$-Higgs 丛给出 $(\mathrm{B},\mathrm{A},\mathrm{A})$-膜；对多数 $p, q \in \N$，这些膜需要同时借助阿贝尔数据与非阿贝尔数据来描述其在 Hitchin 纤维化中的位置。由定义 2.13，一般 $\SO(p+q,p)$-Higgs 丛 $(V\oplus W,\Phi)$ 的特征多项式可写成
\[
    \det(\Phi-\eta) = \eta^q(\eta^{2p} + a_1x^{2p-2} + \cdots + a_p),
\]
并定义由下式给出的、一般光滑的谱曲线 $\pi : S \to \Sigma$：
\begin{equation}
    S = \{0 = \eta^{2p} + a_1x^{2p-2} + \cdots + a_p\}, \tag{4.13}
\end{equation}
其中 $a_i \in H^0(\Sigma, K^{2i})$，$\eta$ 为 $\pi^*K$ 的 重言截面。与其他情形一样，由于 (4.13) 中定义 $S$ 的多项式只含偶次项，存在自然对合 $\sigma : \eta \mapsto -\eta$，由此可定义商曲线 $\bar\pi : \bar S := S/\sigma \to \Sigma$。

由 [11, Theorem 4.11]，$\SO(p+q,p)$-Higgs 丛与谱数据三元组 $(L,M,\tau)$ 之间存在一一对应，其中：
\begin{enumerate}[label=(\roman*),leftmargin=2.2em]
\item 一个线丛 $L \in \Jac(\bar S)[2]$，即 $\bar S$ 上次数为零的 2-挠线丛构成的空间；
% [Tutor Note: p.34 原文作 "A line line bundle L"，"line" 重复一次，为笔误；译文只译作"线丛"。]
\item 二重辅助曲线 $\cC := \{0 = \eta^2 - a_p\} \subset \operatorname{Tot}(K^p)$ 上一个型为 $(q-1,1)$ 的秩 $q$ 等变正交丛 $M$（其中 $a_p$ 如 (4.13) 所示），带有一个定向的选取，并满足如下半稳定性条件：所有不变的各向同性子丛 $M' \subset M$ 都满足 $\deg(M') \leqslant 0$；
\item 在 $a_1$ 的每个零点 $x$ 处选取一个等距 $\tau$，把 $W$ 的相应纤维与 $M$ 的 $-1$ 特征空间 $M^-$ 的相应纤维等同起来；这给出了 (i) 与 (ii) 之间所需的扩张数据。该扩张数据可等同于：对 $\cC$ 的每个分歧点 $x$ 指定一个单位向量 $\tau_x \in L_x$，即一族 $\{\tau_x\}$。
\end{enumerate}

在上述描述中，谱曲线 $\cC$ 带有换叶对合（sheet swapping involution）$\sigma_{\cC} : \cC \to \cC$，其定义为 $\sigma_{\cC}(\eta) = -\eta$。于是，若该对合到 $M$ 的提升具有维数分别为 $q-1$ 与 $1$ 的 $\pm1$ 特征空间，则称等变正交丛 $(M,Q_M,\tilde\sigma_{\cC})$ 的型为 $(q-1,1)$。
